\documentclass[11pt,a4paper]{article}
\usepackage[T2A]{fontenc}
\usepackage[utf8]{inputenc}
\usepackage[russian]{babel}
\usepackage[margin=19mm,headheight=15pt]{geometry}
\usepackage{amsmath,amssymb,booktabs,array,xcolor,enumitem,fancyhdr,graphicx,placeins}
\usepackage{xurl}
\usepackage[colorlinks=true,linkcolor=blue!50!black,urlcolor=blue!50!black]{hyperref}
\setlength{\parindent}{0pt}
\setlength{\parskip}{5pt}
\setlength{\emergencystretch}{3em}
\setlist[itemize]{leftmargin=*,topsep=3pt,itemsep=2pt,parsep=0pt}
\setlist[enumerate]{leftmargin=*,topsep=3pt,itemsep=2pt,parsep=0pt}
\pagestyle{fancy}
\fancyhf{}
\fancyhead[L]{EM и гауссовские смеси}
\fancyhead[R]{Конспект преподавателя}
\fancyfoot[C]{\thepage}
\newcommand{\E}{\mathbb E}
\newcommand{\R}{\mathbb R}
\newcommand{\Normal}{\mathcal N}
\newcommand{\KL}{D_{\mathrm{KL}}}
\newcommand{\LSE}{\operatorname{logsumexp}}
\newcommand{\diag}{\operatorname{diag}}
\newcommand{\tr}{\operatorname{tr}}
\newcommand{\board}[1]{\par\textbf{На доске.} #1\par}
\newcommand{\teach}[1]{\par\textbf{Комментарий преподавателю.} #1\par}
\begin{document}
\begin{center}
 {\LARGE\bfseries EM алгоритм.\\[3pt] Модель гауссовской смеси.}\\[7pt]
 {\large Конспект преподавателя и разбор заданий}\\[4pt]
 ИППИ РАН \textbar\ Байесовские методы машинного обучения \textbar\ Осень 2026
\end{center}

Цель семинара --- понять EM как алгоритм обучения модели со скрытыми
переменными, реализовать гауссовскую смесь и использовать эти идеи в задаче,
где один объект --- целое распределение. Распределение наблюдается через
гистограмму с $B$ бинами. Студенты должны сами определить полезное сходство
распределений, предложить признаки или вероятностную модель и проверить
свой выбор. В этой задаче существенны форма, порядок бинов и точность
эмпирической гистограммы.

Материалы продолжают прошлый семинар про вариационный вывод и VAE.
Используется та же ELBO, но для конечной смеси E-шаг вычисляется точно.
Основной маршрут не требует нейросетей. Ноутбук содержит определения,
задания на реализацию, аналитические проверки, контролируемые опыты и
вопросы к интерпретации; заполненная версия сохраняет реальный запуск.

\section{Проведение семинара}

Предполагаются формула Байеса, ожидание, нормальная плотность, элементарное
дифференцирование и NumPy. Все данные генерируются локально.
Для основных заданий нужны NumPy, SciPy, Matplotlib, scikit-learn и IPython;
PyTorch используется только в приложении. GPU и загрузка датасета не нужны.

\begin{center}
\begin{tabular}{p{.19\linewidth}p{.63\linewidth}r}
\toprule
Этап & Содержание и ожидаемый результат & Минуты\\
\midrule
Модель и вывод & Скрытая метка, наблюдаемое правдоподобие, ELBO, E/M & 20\\
Задания 1--2 & Логарифмы плотностей, мягкие веса, M-шаг, цикл EM & 35\\
Задание 3 & Перенос EM на смесь мультиномиальных распределений & 15\\
Основная задача & Собственный план, признаки, модели, контроль и обсуждение & 50\\
\bottomrule
\end{tabular}
\end{center}

Приложения про транспорт, выбор $K$ и VAE предназначены для продолжения
или самостоятельной работы. Для 90-минутного занятия можно предоставить
готовый цикл EM и обсуждать мультиномиальный M-шаг без его реализации;
основную задачу сохранить как центральную часть. Перед открытием ноутбука
с решениями следует дать студентам пять минут на собственные идеи.

\teach{В студенческой версии незаполненные функции вызывают
\texttt{NotImplementedError}. Только это исключение переводит проверку
в состояние <<задание ещё не выполнено>>; настоящие ошибки не скрываются.
После заполнения нужно повторно выполнить определение функции, проверку
и зависимые эксперименты. Итоговый ARI не является условием прохождения теста.}

\section{Гауссовская смесь и скрытые переменные}

Пусть наблюдения $x_1,\ldots,x_N\in\R^d$ независимы. Вводим
скрытые метки $z_i\in\{1,\ldots,K\}$:
\[
 p(z_i=k)=\pi_k,\qquad
 p(x_i\mid z_i=k)=\Normal(x_i;\mu_k,\Sigma_k),\qquad
 \pi_k>0,\quad\sum_k\pi_k=1,\quad\Sigma_k\succ0.
\]
Нормальная плотность равна
\[
 \Normal(x;\mu,\Sigma)=
 (2\pi)^{-d/2}|\Sigma|^{-1/2}
 \exp\!\left[-\tfrac12(x-\mu)^\top\Sigma^{-1}(x-\mu)\right].
\]
Совместная модель и маргинальная плотность имеют вид
\[
 p_\theta(x_i,z_i=k)=\pi_k\Normal(x_i;\mu_k,\Sigma_k),\qquad
 p_\theta(x_i)=\sum_k\pi_k\Normal(x_i;\mu_k,\Sigma_k).
\]
Параметры $\theta=\{\pi_k,\mu_k,\Sigma_k\}_{k=1}^K$ оцениваем по
наблюдаемому логарифму правдоподобия
\[
 \ell(\theta)=\sum_i\log\sum_k\pi_k\Normal(x_i;\mu_k,\Sigma_k).
\]
Это сумма логарифмов плотностей, а не вероятностей отдельных точек.
Поэтому значение $\ell$ может быть положительным. Логарифмы натуральные.

Если бы метки были известны, полное логарифмическое правдоподобие было бы
\[
 \log p_\theta(X,z)=\sum_{i,k}\mathbf1\{z_i=k\}
 [\log\pi_k+\log\Normal(x_i;\mu_k,\Sigma_k)].
\]
Каждая компонента обучалась бы отдельно. EM заменяет неизвестные индикаторы
их условными ожиданиями, сохраняя неопределённость назначения.

\board{Чем отличаются $\pi_k$ и $p(z_i=k\mid x_i)$? Первый вес описывает
долю компоненты до наблюдения объекта; второй зависит от конкретного $x_i$.
Перестановка всех номеров компонент не меняет плотность. Поэтому номер
кластера сам по себе не имеет смысла. Число компонент также не обязано
совпадать с числом мод: близкие гауссианы могут давать одну моду.}

\teach{В этом семинаре EM оценивает параметры точечно методом ML или MAP.
Точный постериор E-шага относится к $z_i$ при фиксированном $\theta$.
Это не байесовский постериор по параметрам смеси и не алгоритм их сэмплирования.}

\section{От ELBO к E- и M-шагам}

Для каждого объекта введём произвольное распределение $q_i(k)$.
Определим
\[
 \mathcal F(q,\theta)=\sum_{i,k}q_i(k)
 [\log p_\theta(x_i,z_i=k)-\log q_i(k)].
\]
При соглашении $0\log0=0$ и корректном покрытии носителя раскрываем KL:
\begin{align*}
 \KL(q_i\|p_\theta(z_i\mid x_i))
 &=\sum_kq_i(k)\log q_i(k)
   -\sum_kq_i(k)\log p_\theta(x_i,z_i=k)+\log p_\theta(x_i),\\
 \ell(\theta)&=\mathcal F(q,\theta)
       +\sum_i\KL(q_i\|p_\theta(z_i\mid x_i)).
\end{align*}
Неотрицательность KL даёт нижнюю границу. Параметр $q$ нужен для её
оптимизации, а не как дополнительная часть генеративной модели.

\subsection{E-шаг}
При текущем $\theta^{(t)}$ максимизируем границу по $q$.
Оптимум достигается на точном постериоре:
\[
 r_{ik}=q_i^{(t)}(k)=
 \frac{\pi_k^{(t)}\Normal(x_i;\mu_k^{(t)},\Sigma_k^{(t)})}
 {\sum_l\pi_l^{(t)}\Normal(x_i;\mu_l^{(t)},\Sigma_l^{(t)})}.
\]
Для каждого $i$ веса неотрицательны и суммируются в единицу.
Число $r_{ik}$ называют ответственностью компоненты за объект.
Жёсткое назначение $\arg\max_k r_{ik}$ удобно для рисунка, но в M-шаг
передаются сами мягкие веса.

\subsection{M-шаг и гарантия подъёма}
При фиксированном $q^{(t)}$ энтропия не зависит от новых параметров.
Поэтому максимизируем
\[
 Q(\theta\mid\theta^{(t)})=
 \sum_{i,k}r_{ik}\log p_\theta(x_i,z_i=k).
\]
При точной максимизации или хотя бы увеличении этой функции
\[
 \ell(\theta^{(t+1)})
 \geq\mathcal F(q^{(t)},\theta^{(t+1)})
 \geq\mathcal F(q^{(t)},\theta^{(t)})
 =\ell(\theta^{(t)}).
\]
Равенство справа обеспечено E-шагом. После изменения $\theta$ зазор KL
обычно снова положителен, пока не выполнен следующий E-шаг.
Монотонность не означает ни глобального оптимума, ни существования
конечного максимума, ни полезности получившихся кластеров.

\board{Предложите студентам восстановить цепочку самостоятельно.
Типичная ошибка --- считать, что граница совпадает с правдоподобием
для любых $q$ и $\theta$. Другая ошибка --- перепутать условное ожидание
логарифма совместной модели с логарифмом ожидаемого правдоподобия.}

\subsection{Связь с VAE}
Вариационный вывод ограничивает семейство $q$, например нейросетевым
гауссовским энкодером. Тогда точный постериор может не принадлежать этому
семейству, и KL-зазор не обязан обращаться в ноль. В конечной смеси
E-шаг --- точное суммирование по $K$ меткам, а не Монте-Карло и не
репараметризация дискретной случайной величины. Нейросетевой M-шаг обычно
также заменяется градиентными обновлениями.

\section{Вывод M-шага для гауссовской смеси}

Положим $N_k=\sum_i r_{ik}$. Часть $Q$, зависящая от весов,
равна $\sum_kN_k\log\pi_k$. С множителем Лагранжа для $\sum_k\pi_k=1$
получаем $N_k/\pi_k-\lambda=0$, откуда
\[
 \pi_k^{\rm new}=N_k/N.
\]
При фиксированной ковариации производная по среднему равна
$\Sigma_k^{-1}\sum_i r_{ik}(x_i-\mu_k)$, поэтому
\[
 \mu_k^{\rm new}=\frac1{N_k}\sum_i r_{ik}x_i.
\]
Теперь центрируем по \emph{новому} среднему и обозначим
$S_k=N_k^{-1}\sum_i r_{ik}(x_i-\mu_k^{\rm new})(x_i-\mu_k^{\rm new})^\top$.
Часть цели для ковариации:
\[
 Q_k=-\frac{N_k}{2}\bigl[\log|\Sigma_k|+
                         \tr(\Sigma_k^{-1}S_k)\bigr]+\mathrm{const}.
\]
При положительно определённом $S_k$ оптимум равен
\[
 \Sigma_k^{\rm new}=S_k.
\]
Если $S_k$ вырожденна, без дополнительных ограничений внутреннего
положительно определённого максимума может не быть. Деление производится
на $N_k$, а не на $N_k-1$: это ML-обновление, не несмещённая оценка дисперсии.

\subsection{Диагональная модель и фиксированный нижний порог}
В собственной реализации $\Sigma_k=\diag(v_{k1},\ldots,v_{kd})$.
Взвешенная покоординатная дисперсия равна
\[
 s_{kj}^2=\frac1{N_k}\sum_i r_{ik}(x_{ij}-\mu_{kj}^{\rm new})^2.
\]
Зафиксируем $v_{kj}\geq v_{\min}>0$. Для одной координаты цель имеет вид
$-N_k[\log v+s^2/v]/2$, её производная пропорциональна $(s^2-v)/v^2$.
Следовательно, точный M-шаг в этом ограниченном семействе:
\[
 v_{kj}^{\rm new}=\max\{s_{kj}^2,v_{\min}\}.
\]
Такой шаг сохраняет доказательство подъёма ограниченной задачи.
Операция $s^2+\varepsilon$ отличается от $\max(s^2,v_{\min})$.
В \texttt{sklearn} параметр \texttt{reg\_covar} добавляет величину к диагонали;
поэтому буквального совпадения двух реализаций не требуется.

\subsection{Вырождение, старты и критерий остановки}
Без ограничений одна компонента может стянуть ковариацию к нулю около
отдельной точки, делая плотность там неограниченной. Рост обучающего
правдоподобия тогда сопровождается плохой моделью данных.
В семинаре используем нижний порог дисперсий, несколько стартов и проверку
отложенной плотности там, где сравниваются одинаковые наблюдения.

Если $N_k<10^{-10}$, данный запуск прекращается. Случайное восстановление
пустой компоненты --- возможная практическая эвристика, но оно уже не
является тем точным M-шагом, для которого выведена монотонность.
Истории разных запусков нельзя склеивать в одну кривую подъёма.

В коде история содержит начальную и финальную цели. Остановка происходит,
когда улучшение не превосходит
$\texttt{tol}\,(1+|\ell_{\rm previous}|)$; дополнительно установлен предел
итераций. Из независимых стартов выбирается лучший по обучающей цели.
Флаг достижения критерия и число итераций нужно сообщать вместе с результатом.

\section{Разбор заданий 1--2 и игрушечных опытов}

\subsection{Стабильный E-шаг и формы массивов}
Для диагональной модели
\[
 \log\Normal(x_i;\mu_k,\diag v_k)
 =-\frac12\sum_j\left[\log(2\pi v_{kj})+
                           \frac{(x_{ij}-\mu_{kj})^2}{v_{kj}}\right].
\]
Вычисляем $a_{ik}=\log\pi_k+\log\Normal(x_i;\mu_k,\diag v_k)$ и
\[
 s_i=\LSE_k(a_{ik}),\qquad r_{ik}=\exp(a_{ik}-s_i),\qquad
 \ell=\sum_i s_i.
\]
Тождество $\LSE(a)=m+\log\sum_k e^{a_k-m}$ с $m=\max_k a_k$
объясняет устойчивость. Прямое вычисление маленьких плотностей может дать
нули, а последующая нормировка --- неопределённость $0/0$.

\begin{center}
\begin{tabular}{lll}
\toprule
Величина & Форма массива & Смысл осей\\
\midrule
$X$ & $(N,d)$ & объект, координата\\
$\mu,v$ & $(K,d)$ & компонента, координата\\
$a,r$ & $(N,K)$ & объект, компонента\\
$\pi,N_k$ & $(K,)$ & компонента\\
\bottomrule
\end{tabular}
\end{center}

\board{Для $d=1$, $\pi_1=\pi_2=1/2$, $\mu_1=0$, $\mu_2=2$,
$v_1=v_2=1$ и $x=0$ разность логарифмических оценок равна $a_1-a_2=2$.
Значит $r_1=1/(1+e^{-2})\approx0.8808$. Для $x=1$ обе ответственности
равны $1/2$. Для $x=10000$ вычисления в логарифмах остаются конечными.}

\subsection{Контроль M-шага}
Для точек $0,2,8,10$ и жёстких групп $\{0,2\}$, $\{8,10\}$
должны получиться веса $(1/2,1/2)$, средние $(1,9)$ и дисперсии $(1,1)$.
Если все точки совпали, дисперсия после ограниченного M-шага равна
$v_{\min}$. Далее проверяются нормировка ответственностей, подъём цели
и совпадение последнего значения истории с правдоподобием возвращённой модели.
Тесты не требуют восстановления конкретной нумерации компонент.

\subsection{Что обсуждать на рисунках}
Первый пример состоит из трёх диагональных гауссиан. На одном рисунке
показаны назначения, энтропия $-\sum_k r_{ik}\log r_{ik}$ и история цели.
Энтропия выше в областях, где компоненты дают сопоставимые плотности.
Это неопределённость скрытой метки при фиксированных оценённых параметрах,
а не оценка неопределённости самих параметров.

\begin{figure}[htbp]
\centering\includegraphics[width=\linewidth]{assets/toy_em.png}
\caption{Игрушечный пример из заполненного ноутбука: мягкие назначения,
неоднозначность компоненты и подъём наблюдаемого логарифма правдоподобия.}
\end{figure}

В контрольном запуске наша диагональная модель и диагональная
\texttt{GaussianMixture} дают средний обучающий $\log p(x)\approx-2.9772$.
Это удобная проверка согласованности, а не гарантия совпадения всех стартов.
Во втором примере компоненты вытянуты и повёрнуты. Для общей отложенной
выборки средний $\log p(x)$ равен примерно $-3.4660$ у диагональной и
$-3.1948$ у полной модели. Полная ковариация описывает наклон облаков.
Здесь значения сравнимы: обе плотности определены в исходном $\R^2$.

\subsection{Связь с K-means}
При фиксированной общей ковариации $\sigma^2I$ и равных весах
\[
 r_{ik}\propto\exp[-\|x_i-\mu_k\|^2/(2\sigma^2)].
\]
При $\sigma\to0$ и единственном ближайшем среднем ответственность
сосредоточивается на нём. M-шаг средних превращается в обновление K-means.
Это объясняет связь, но не делает произвольную GMM тождественной K-means:
разные веса и ковариации меняют границы, а конечная дисперсия даёт мягкие веса.
\FloatBarrier

\section{Основная постановка: объектом является распределение}

Наблюдается скалярная случайная величина $T$, но кластеризовать требуется
не её отдельные реализации, а распределения разных объектов.
У объекта $i$ имеется $n_i$ реализаций $T$, сведённых в общие бины
$[a_{b-1},a_b)$ ширины $\Delta_b=a_b-a_{b-1}>0$.
\[
 c_{ib}=\text{число отсчётов в бине }b,\qquad
 p_{ib}=c_{ib}/n_i,\qquad h_{ib}=p_{ib}/\Delta_b.
\]
Это три разных представления:
\[
 \sum_b c_{ib}=n_i,\qquad\sum_b p_{ib}=1,\qquad
 \sum_b h_{ib}\Delta_b=1.
\]
Если пользователь передаёт высоты плотности $x_i=h_i$, для методов,
работающих с вероятностями бинов, сначала нужно вычислить $p_{ib}=h_{ib}\Delta_b$.
Из нормированной плотности нельзя восстановить $n_i$ без дополнительной информации.

\board{Для ширин $(1,2)$ и высот $(0.5,0.25)$ массы равны $(0.5,0.5)$.
Сумма высот равна $0.75$, а интеграл плотности --- единице.
При слиянии бинов складываются счётчики или массы; новая высота равна
суммарной массе, делённой на новую ширину.}

Матрица признаков имеет форму $(N,B)$: строка --- распределение,
столбец --- бин. Разбиение должно назначить один кластер всей строке.
Если обучить GMM на отдельных скалярных отсчётах, получится другая задача:
моделирование плотности внутри распределения или объединённой выборки.
Перестановка бинов обычно меняет смысл распределения на прямой.
Разные сетки сначала приводятся к общей сетке с сохранением масс;
уточнить расположение массы внутри исходного бина без предположений нельзя.

\teach{Не начинайте основную задачу с готового sqrt-преобразования.
Сначала попросите пары сформулировать, что они считают сходством:
положение массы, число мод, асимметрию, масштаб или форму с точностью до сдвига.
Разные ответы ведут к разным корректным методам. Отсутствует единственная
<<правильная>> геометрия без содержательного определения кластеров.}

\section{Учебные данные и честный протокол}

В ноутбуке генерируются $N=720$ объектов, $B=64$ равных бина на $[-4,4]$,
размеры выборок $n_i$ от 60 до 600. Известно $K=4$ --- это условие учебной
задачи, а не вывод по истинным меткам. Базовые семейства:
\begin{enumerate}
\item стандартная нормальная плотность;
\item симметричная смесь двух нормальных плотностей со средними $\pm0.85$
и стандартным отклонением $0.35$, после деления переменной на
$s=\sqrt{0.85^2+0.35^2}$;
\item смесь $0.8\Normal(-0.45,0.45^2)+0.2\Normal(1.8,0.6^2)$,
после деления переменной на $s=\sqrt{1.044}$;
\item отражение третьего семейства относительно нуля.
\end{enumerate}
У третьей смеси среднее $0.8(-0.45)+0.2(1.8)=0$, второй момент равен
$0.8(0.45^2+0.45^2)+0.2(1.8^2+0.6^2)=1.044$.
Поэтому у всех четырёх базовых форм теоретические среднее и дисперсия равны
0 и 1. При этом число мод, асимметрия и хвосты отличаются.

Далее к каждому объекту независимо применяются сдвиг
$a_i\sim\Normal(0,0.18^2)$ и растяжение
$s_i=\exp(\eta_i)$, $\eta_i\sim\Normal(0,0.1^2)$.
Законы этих преобразований одинаковы для всех семейств. Массы бинов
считаются разностями аналитических CDF, нормируются внутри $[-4,4]$,
затем генерируются мультиномиальные счётчики. Обрезание хвостов, сетка
и конечная выборка нарушают точное равенство эмпирических моментов;
не следует обещать его для реализованных гистограмм.

\begin{figure}[htbp]
\centering\includegraphics[width=\linewidth]{assets/histogram_examples.png}
\caption{Восемь наблюдаемых объектов. Различия формы смешаны с шумом
счётчиков и небольшими индивидуальными сдвигами и растяжениями.}
\end{figure}

Генератор использует \texttt{default\_rng(2026)}. Объекты случайно
разделяются с \texttt{default\_rng(42)}: 480 для обучения, 240 для теста.
Разделяются целые строки; не бины и не части выборки одного объекта.
PCA, стандартизация, GMM, шаблоны, медоиды и энкодер обучаются только на
обучающих объектах. Тестовые гистограммы используются для назначения
новых объектов и заключительной оценки.

Истинные метки и истинные массы доступны через \texttt{simulation\_audit}.
Это учебное сокрытие ответа, а не защита от чтения исходника. Открывать
их нужно после выбора методов и параметров. В реальных данных меток
обычно нет. Основной seed и все параметры фиксируются до аудита.
\FloatBarrier

\section{Три способа задать модель распределений}

\subsection{GMM по признакам}
Простейший вектор --- среднее и дисперсия по центрам бинов $t_b$:
\[
 m_i=\sum_bp_{ib}t_b,\qquad
 v_i=\sum_bp_{ib}(t_b-m_i)^2.
\]
Это моменты дискретной меры в центрах бинов, приближающие моменты исходного
распределения. В данной симуляции такой вектор намеренно теряет основную
информацию о семействе. Полезное расширение --- асимметрия
$\sum_bp_{ib}(t_b-m_i)^3/v_i^{3/2}$ и четвёртый стандартизованный момент
$\sum_bp_{ib}(t_b-m_i)^4/v_i^2$. В коде используется именно четвёртый
момент, без вычитания трёх. Он чувствителен к редким отсчётам в хвостах.

Другой вариант --- вектор всех масс. Он лежит на симплексе
$\sum_bp_b=1$, поэтому полная ковариация в $\R^B$ вырождена.
Регуляризация или снижение размерности помогают численно, но вводят
дополнительные модельные предположения. GMM в признаках моделирует
распределение векторов признаков, а не обязательно корректную плотность
на исходном пространстве распределений.

\subsection{Преобразование Хеллингера и PCA}
Для дискретных вероятностей используем соглашение
\[
 H^2(p,q)=\frac12\sum_b(\sqrt{p_b}-\sqrt{q_b})^2.
\]
Поэтому $p\mapsto\sqrt p$ переводит эту геометрию в евклидову с постоянным
множителем $\sqrt2$ для расстояния. Нулевые массы допустимы, сглаживание
для вычисления квадратного корня не требуется. Для неравных бинов сначала
переходим от высот к массам; применение корня к одним высотам неверно.

В эталонном подходе: $\sqrt p\to\mathrm{PCA}(8)\to\mathrm{GMM}$
с полной ковариацией, шестью стартами и \texttt{reg\_covar=1e-4}.
PCA обучается только на training; whitening не используется.
Ортогональная проекция с отброшенными координатами теряет часть расстояния.
Этот выбор --- разумная отправная точка, а не единственное решение.

\subsection{Смесь мультиномиальных распределений}
Если известны счётчики, можно непосредственно моделировать точность
гистограммы. У \emph{всей строки} один скрытый номер $z_i$:
\[
 z_i\sim\operatorname{Cat}(\pi),\qquad
 c_i\mid z_i=k,n_i\sim\operatorname{Multinomial}(n_i,\theta_k),
 \qquad\sum_b\theta_{kb}=1.
\]
Размер $n_i$ считается заданным, а не случайной переменной, предсказывающей
класс. Условная вероятность равна
\[
 p(c_i\mid k,n_i)=\frac{n_i!}{\prod_bc_{ib}!}\prod_b\theta_{kb}^{c_{ib}}.
\]
Мультиномиальный коэффициент не зависит от $k$, поэтому E-шаг:
\[
 a_{ik}=\log\pi_k+\sum_bc_{ib}\log\theta_{kb},\qquad
 r_{ik}=\exp(a_{ik}-\LSE_l a_{il}).
\]
При ML M-шаг даёт
\[
 \pi_k=\frac{N_k}{N},\qquad
 \theta_{kb}=\frac{\sum_ir_{ik}c_{ib}}{\sum_ir_{ik}n_i}.
\]
Для вывода достаточно максимизировать $\sum_b A_{kb}\log\theta_{kb}$
при $\sum_b\theta_{kb}=1$, где $A_{kb}=\sum_ir_{ik}c_{ib}$.
Нормировка производится по взвешенному количеству отсчётов, а не по $N_k$.

\subsection{MAP и разбор задания 3}
Чтобы избежать нулевых шаблонов, в реализации задан независимый prior
$\theta_k\sim\operatorname{Dirichlet}(\alpha,\ldots,\alpha)$,
$\alpha=1.05>1$. Его логарифм добавляет
$(\alpha-1)\sum_{k,b}\log\theta_{kb}$. Отсюда MAP-обновление:
\[
 \theta_{kb}^{\rm new}=
 \frac{\sum_ir_{ik}c_{ib}+\alpha-1}
      {\sum_ir_{ik}n_i+B(\alpha-1)}.
\]
Именно $\alpha-1$, а не $\alpha$, соответствует максимуму апостериорной
плотности. Для известного назначения апостериорное среднее использовало
бы другую добавку; это другая оценка. В коде требуется $\alpha>1$.

История содержит
\[
 J(\pi,\theta)=\sum_i\LSE_k\left(\log\pi_k+
                       \sum_bc_{ib}\log\theta_{kb}\right)
             +(\alpha-1)\sum_{k,b}\log\theta_{kb}.
\]
Она отличается от log posterior только константами для фиксированных
данных и $\alpha$. Монотонность проверяется для $J$, а не для чистого
логарифма правдоподобия. Для абсолютного log probability нужно вернуть
$\log(n_i!)-\sum_b\log(c_{ib}!)$, устойчиво вычисляя их через
\texttt{gammaln}. В назначениях этот член сокращается.

\board{Для счётчиков $(8,2,0)$ и $(4,1,0)$, полностью принадлежащих одной
компоненте, взвешенная сумма равна $(12,3,0)$. При $\alpha=2$ новый шаблон
равен $(13,4,1)/18$. Этот пример проверяет добавку и нормировку.
Гистограммы $(9,1,0)$ и $(90,10,0)$ имеют одинаковые доли, но вторая даёт
более концентрированное назначение между различающимися шаблонами.}

\teach{Подстановка долей вместо счётчиков заменяет выбранную модель:
логарифмический вклад объекта перестаёт масштабироваться с $n_i$.
Большие $n_i$ намеренно дают больше информации о шаблоне. Однако один
шаблон не описывает индивидуальные сдвиги и растяжения; при неверной
модели большое $n_i$ может приводить к чрезмерно уверенным назначениям.
Для более богатой модели можно обсуждать случайные параметры формы
или смесь распределений Дирихле--мультиномиальных счётчиков.}

\section{Открытое задание и организация обсуждения}

До демонстрации подсказок попросите студентов письменно ответить:
\begin{enumerate}
\item Какие свойства распределения определяют кластер и какие считаются помехой?
\item Какие признаки или расстояния сохраняют эти свойства?
\item Что будет потеряно при таком представлении?
\item Как учитывать различное $n_i$ и наличие нулевых бинов?
\item Как назначить кластер новому объекту без переобучения на тесте?
\end{enumerate}
Функция \texttt{fit\_distribution\_clustering(P\_fit,P\_new,K,seed)}
возвращает метки обучающих и новых объектов и текстовое описание метода.
Первый подход должен использовать GMM. Второй может менять признаки,
геометрию или условную модель; например, четыре момента, смесь счётчиков,
транспортные медоиды или VAE-представление.

Далее требуются изображения нескольких отдельных гистограмм и средней
плотности каждого кластера, численность групп и распределение $n_i$.
Среднее по объектам не тождественно мультиномиальному шаблону:
шаблон использует мягкие назначения и взвешивает бины количеством отсчётов.
Если кластеры различаются в основном по шумности, это нужно заметить.

Устойчивость оценивается по трём инициализациям. ARI между предсказанными
разбиениями одних и тех же объектов не требует истинных меток.
Но стабильность не доказывает содержательной правильности: устойчивым
может оказаться разбиение по нерелевантному признаку.

В контрольном опыте каждый объект прореживается до 30 отсчётов выбором
без возвращения из имеющихся счётчиков. Все исходные $n_i\geq60$.
Сохраняется раздельность объектов training/test, методы переобучаются
по тому же заранее заданному протоколу. Это зависимое представление тех
же объектов с меньшей точностью, а не новая независимая выборка объектов.
Другой допустимый контроль --- новая независимая симуляция с иным $B$
или большим разбросом сдвигов. Для другой сетки преобразования переобучаются.

\teach{Удобный порядок обсуждения: сначала ожидание студента, затем график,
затем численные результаты. Попросите объяснить хотя бы один случай,
где выбранный метод должен ошибаться. Не сводите работу к поиску лучшего ARI.}

\section{Оценка и разбор контрольного запуска}

Истинные метки открываются после фиксации методов. ARI учитывает совпадение
пар объектов и корректирует случайное совпадение; значение 1 означает
одинаковое разбиение, значения около нуля соответствуют случайному уровню.
NMI --- нормированная взаимная информация в диапазоне от 0 до 1.
Обе метрики инвариантны к перестановке номеров кластеров.
Обычная accuracy без сопоставления номеров здесь неинформативна.

\begin{figure}[htbp]
\centering\includegraphics[width=\linewidth]{assets/shape_and_moments.png}
\caption{После открытия симуляции: средние истинные формы заметно отличаются,
а признаки из первых двух моментов сильно перекрываются.}
\end{figure}

\begin{center}
\begin{tabular}{lrr}
\toprule
Метод & Test ARI & Test NMI\\
\midrule
GMM: среднее и дисперсия & $-0.0026$ & $0.0115$\\
GMM: четыре момента & $0.9903$ & $0.9850$\\
$\sqrt p$ + PCA(8) + GMM & $0.9889$ & $0.9846$\\
EM: мультиномиальная смесь & $0.9403$ & $0.9235$\\
$W_1$: медоиды & $0.7696$ & $0.7669$\\
VAE(3) + GMM, приложение & $0.9467$ & $0.9377$\\
\bottomrule
\end{tabular}
\end{center}

Это сохранённый контрольный запуск: Python 3.9.7, NumPy 1.19.5,
SciPy 1.10.1, scikit-learn 0.24.2, PyTorch 2.4.1, CPU.
Числа приведены для интерпретации, а не как пороги оценивания.
Версии библиотек и случайные старты могут немного менять результат.
Первые два момента не сохраняют различия базовых форм. Четыре момента
достаточны для данного специально сконструированного семейства, но в общем
случае не определяют распределение. PCA по корням масс сохраняет более
богатую информацию о форме. Счётная смесь успешно использует распределение
по бинам, хотя её фиксированные шаблоны приближают внутрикластерную вариативность.

При прореживании до $n_i=30$ ARI равен примерно $-0.0033$ для двух моментов,
$0.6991$ для sqrt/PCA/GMM и $0.8226$ для мультиномиальной смеси.
Уменьшение $n_i$ ухудшает оценку формы; в этом опыте модель счётчиков
сохраняет больше информации о семьях. Нельзя переносить ранжирование
на произвольные данные. Три старта GMM в выбранных признаках дают
одинаковые отложенные разбиения, но это не устраняет модельную ошибку.

\subsection{Что допустимо сравнивать по правдоподобию}
Логарифмы плотностей в разных пространствах признаков имеют разные меры
и размерности. Значения GMM по моментам, GMM после PCA и мультиномиального
log probability нельзя непосредственно ранжировать между собой.
Потерянная при PCA информация не возвращается простым сравнением оценок
плотности в проекции. Сравнение diagonal/full на одном исходном $x$
корректно; межмодельный аудит кластеров в основной задаче выполняется
общими метриками разбиения и интерпретацией.
\FloatBarrier

\section{Приложение: порядок бинов и расстояние Васерштейна}

Для вероятностных мер на общей прямой
$W_1(p,q)=\int|F_p(t)-F_q(t)|\,dt$. Если все массы сосредоточены в
общих упорядоченных центрах $t_1<\cdots<t_B$, CDF постоянны между центрами:
\[
 W_1(p,q)=\sum_{b=1}^{B-1}
 \left|\sum_{j\leq b}(p_j-q_j)\right|(t_{b+1}-t_b).
\]
Это точная формула для дискретных мер в центрах. Для непрерывной плотности,
постоянной внутри бинов, CDF внутри бина линейна и интегрирование модуля
её разности требует отдельного расчёта. Эти две интерпретации не следует
смешивать. Реализация проверяется на расстояниях между точечными массами
и сопоставляется с \texttt{scipy.stats.wasserstein\_distance}.

Wasserstein учитывает, насколько далеко нужно перенести массу.
Хеллингер между точечной массой в бине $b$ и массой в другом бине одинаков
независимо от их расстояния; $W_1$ растёт с физическим расстоянием центров.
Поэтому транспортная геометрия естественна для положения массы, но может
сильнее реагировать на сдвиги, которые мы хотели считать несущественными.

Обычное среднее минимизирует сумму квадратов евклидовых расстояний;
замена расстояния в K-means сама по себе не даёт нужного обновления центра.
В приложении применяются медоиды:
\begin{enumerate}
\item выбрать $K$ обучающих объектов-представителей;
\item назначить каждому объекту ближайший медоид;
\item внутри каждой группы выбрать объект с минимальной суммой расстояний
до остальных членов;
\item повторять до стабилизации, использовать несколько стартов.
\end{enumerate}
Такое чередование даёт локальное решение. Матрица попарных расстояний
требует памяти $O(N^2)$. Новые объекты сравниваются с обучающими медоидами,
а не включаются в их переоценку. Более низкий ARI транспорта в этом запуске
объясним несовпадением геометрии с заданными семействами и локальной оптимизацией;
это не свидетельство некорректности расстояния.

\section{Приложение: неизвестное число компонент}

В основной задаче $K$ известно. В другой задаче используйте внутреннюю
валидацию, устойчивость, интерпретируемость и, с оговорками, BIC:
\[
 \mathrm{BIC}=-2\ell(\widehat\theta)+m\log N,\qquad
 m=(K-1)+Kd+K\frac{d(d+1)}2
\]
для полной GMM в фиксированном $d$-мерном пространстве. Меньше --- лучше.
В ноутбуке сравниваются $K=1,\ldots,7$ в одной фиксированной PCA-проекции.
Выбор пространства и его обучение не меняются от значения $K$.
Смеси имеют особые точки и перестановочную неидентифицируемость;
обычная регулярная асимптотическая аргументация BIC здесь не является
универсальной гарантией. Компонент может понадобиться больше, чем
содержательных групп, чтобы описать негауссовскую форму одного семейства.
Тестовый ARI не используется для выбора $K$.

\section{Приложение: VAE для гистограмм и последующая GMM}

Обсудить этот вариант полезно после прошлого семинара. Генеративная модель:
\[
 z_i\sim\Normal(0,I_J),\qquad
 c_i\mid z_i,n_i\sim\operatorname{Multinomial}(n_i,\rho_\theta(z_i)),
 \quad\rho_\theta(z)=\operatorname{softmax}(g_\theta(z)).
\]
Энкодер задаёт
$q_\phi(z_i\mid c_i,n_i)=\Normal(\mu_{\phi,i},\diag(e^{s_{\phi,i}}))$.
Он получает $\sqrt{p_i}$ и стандартизованный $\log n_i$;
параметры этой стандартизации оцениваются только на training.
Размер выборки помогает различать точность одинаковых эмпирических долей.

Отрицательный ELBO для одного объекта, без не зависящего от параметров
мультиномиального коэффициента, равен
\[
 -\mathcal L_i=
 -\E_q\sum_bc_{ib}\log\rho_{\theta b}(z_i)
 +\frac12\sum_{j=1}^J\left(\mu_{ij}^2+e^{s_{ij}}-1-s_{ij}\right)
 +\mathrm{const}(c_i).
\]
В обучении используется один репараметризованный сэмпл и коэффициент KL,
равный единице. Деление реконструкции на $n_i$ изменяет соотношение с KL
и перестаёт быть записанным ELBO. Бинарная cross-entropy по нормированным
бинам не соответствует выбранной мультиномиальной модели.

В коде $J=3$, энкодер имеет скрытые слои 48 и 32 с tanh, декодер --- 32
и 48; на выходе логиты для softmax. Обучение: Adam, шаг $2\cdot10^{-3}$,
батч 80, 180 заранее фиксированных эпох на CPU. Логарифмическая дисперсия
энкодера ограничена для устойчивости. Тест не выбирает эпоху или чекпойнт.
Стохастическая история обучения VAE не обязана монотонно убывать как
точная цель EM.

Затем энкодер фиксируется, GMM обучается по средним $\mu_\phi$ обучающих
объектов и назначает кластеры отложенным средним. Это двухэтапная эвристика:
кластерная переменная не входила в исходный VAE. Нормальная априорная
плотность не обещает $K$ разделённых групп, а хорошая реконструкция
не гарантирует хорошую кластеризацию. Для совместной модели потребовалось
бы явно ввести дискретную компоненту и согласовать генеративную модель,
вариационное семейство и целевую функцию.

Декодер при $\mu_\phi$ показывает условные вероятности бинов в одной
латентной точке. Это не интеграл по $q_\phi(z\mid c)$ и не новый случайный
вектор счётчиков. Для генерации гистограммы нужно выбрать $z$, вычислить
вероятности и сэмплировать мультиномиальные счётчики при заданном $n$.
Веса нейросети остаются точечными оценками; вариационный постериор по $z$
не описывает неопределённость этих весов.

\section{Проверки, типичные ошибки и оценивание}

\begin{center}
\begin{tabular}{p{.34\linewidth}p{.59\linewidth}}
\toprule
Наблюдаемый симптом & Что проверить и как объяснить\\
\midrule
Нули или NaN в E-шаге & Вычисление log density, положительность дисперсий,
\texttt{logsumexp}, нормировку по компонентам.\\[3pt]
Цель EM убывает & Старые/новые параметры в истории, центрирование по новому
среднему, мягкие веса; для MAP --- наличие log prior в цели.\\[3pt]
Разные номера правильных групп & Перестановка компонент не меняет модель;
сравнивать разбиения инвариантной метрикой.\\[3pt]
Кластеры по шумности & Показать $n_i$ внутри групп; проверить признаки,
сглаживание и ошибку условной модели.\\[3pt]
Плохая GMM на массах & Проверить симплекс, размерность, вырождение ковариации,
число объектов на компоненту и необходимость проекции.\\[3pt]
Слишком хороший тест & Проверить fit преобразований на training,
момент открытия меток и отсутствие настройки по test ARI.\\[3pt]
VAE хорошо восстанавливает, но плохо группирует & Реконструкция и кластеры
решают разные задачи; показать форму латентного представления и влияние $n_i$.\\
\bottomrule
\end{tabular}
\end{center}

Предлагаемая шкала на 10 баллов: E-шаг и его объяснение --- 2; M-шаг,
цикл EM и монотонность --- 2; мультиномиальная модель и MAP --- 2;
собственные методы основной задачи и честный контроль --- 3;
интерпретация и ограничения --- 1. Приложения могут давать дополнительные
баллы или использоваться в обсуждении. Конкретный уровень ARI в шкалу
не входит. Корректный анализ неудачного метода содержательнее случайно
полученного высокого числа без объяснения.

Итоговые вопросы: почему E-шаг делает границу точной; что оптимизирует
M-шаг; почему несколько стартов не гарантируют глобального решения;
какое ограничение устраняет сжатие дисперсии; почему массы и высоты нельзя
смешивать; чем метка целой гистограммы отличается от метки отсчёта;
когда счётчики дают информацию сверх долей; какие преобразования делают
распределения похожими; что означает успех на синтетическом генераторе.

\clearpage
\section*{Источники и границы материала}
\begin{enumerate}
\item A. P. Dempster, N. M. Laird, D. B. Rubin (1977).
\emph{Maximum Likelihood from Incomplete Data via the EM Algorithm}.
\url{https://doi.org/10.1111/j.2517-6161.1977.tb01600.x}.
\item R. M. Neal, G. E. Hinton (1998).
\emph{A View of the EM Algorithm that Justifies Incremental, Sparse,
and Other Variants}. Авторский PDF:
\url{https://www.cs.toronto.edu/~hinton/absps/emk.pdf}.
\item scikit-learn. \emph{Gaussian mixture models}; описание параметров
\texttt{GaussianMixture}:
\url{https://scikit-learn.org/stable/modules/mixture.html};
\url{https://scikit-learn.org/stable/modules/generated/sklearn.mixture.GaussianMixture.html}.
\item SciPy. \emph{wasserstein\_distance}:
\url{https://docs.scipy.org/doc/scipy/reference/generated/scipy.stats.wasserstein_distance.html}.
\item D. P. Kingma, M. Welling. \emph{Auto-Encoding Variational Bayes}.
\url{https://arxiv.org/abs/1312.6114}.
\end{enumerate}
Генератор распределений, численные сравнения и VAE для счётчиков являются
учебными конструкциями этого семинара, а не воспроизведением экспериментов
перечисленных статей. Основные формулы EM и его вариационная интерпретация
соответствуют классической постановке; практические регуляризации и
модели признаков оговорены отдельно.
\end{document}
